%%%%%%%%%%%%%%%%%%%%%%%%%%%%%%%%%%%%%%%%%%%%%%%%%%%%%%%%%%%%%%%%%%%%%%%%%%%
%%  abbreviations.tex  –  Список сокращений
%%  v9 – Обновлено для структуры из 5 глав с динамическими ГНС,
%%       PQC, платформой, SOC, LLM и предметно-ориентированными терминами.
%%%%%%%%%%%%%%%%%%%%%%%%%%%%%%%%%%%%%%%%%%%%%%%%%%%%%%%%%%%%%%%%%%%%%%%%%%%

\chapter*{Список сокращений}
\addcontentsline{toc}{chapter}{Список сокращений}

\begin{longtable}{@{}p{3.8cm} p{10cm}@{}}
AI (ИИ) & Искусственный интеллект \\
AML & Состязательное машинное обучение \\
APT & Развитая постоянная угроза (Advanced Persistent Threat) \\
AUC-ROC & Площадь под кривой рабочей характеристики приёмника \\
BN & Пакетная нормализация (Batch Normalisation) \\
BWT & Обратный перенос (метрика непрерывного обучения) \\
C\&W & Атака Карлини--Вагнера (Carlini--Wagner Attack) \\
У1, У2, У3 & Обозначения рабочих условий: Состязательная среда, Распределённая динамика, Нестационарность \\
CL-RL & Непрерывное обучение -- Обучение с подкреплением \\
CNN & Свёрточная нейронная сеть (Convolutional Neural Network) \\
CPO & Оптимизация политики с ограничениями (Constrained Policy Optimisation) \\
CT-GNN & Графовая нейронная сеть непрерывного времени \\
CT-TGNN & Темпоральная графовая нейронная сеть непрерывного времени (Метод М1) \\
CVE & Общие уязвимости и экспозиции (Common Vulnerabilities and Exposures) \\
CyberSecLLM & Предметно-ориентированная фундаментальная модель кибербезопасности \\
DDoS & Распределённая атака отказа в обслуживании \\
Dec-CMDP & Децентрализованный марковский процесс принятия решений с ограничениями \\
DNN & Глубокая нейронная сеть (Deep Neural Network) \\
DoS & Атака отказа в обслуживании \\
DP (ДП) & Дифференциальная приватность \\
ECE & Ожидаемая ошибка калибровки (Expected Calibration Error) \\
ELBO & Нижняя граница свидетельства (Evidence Lower Bound) \\
EWC & Упругая консолидация весов (Elastic Weight Consolidation) \\
FedAvg & Федеративное усреднение (Federated Averaging) \\
FedGTD & Федеративная теоретико-игровая защита (Federated Game-Theoretic Defence) \\
FedLLM-API & Федеративная оптимизация графов с большими языковыми моделями и приватностью (Метод М2) \\
FFT (БПФ) & Быстрое преобразование Фурье (Fast Fourier Transform) \\
FGSM & Метод быстрого знака градиента (Fast Gradient Sign Method) \\
FIM & Матрица информации Фишера (Fisher Information Matrix) \\
FWT & Прямой перенос (метрика непрерывного обучения) \\
GAN & Генеративно-состязательная сеть (Generative Adversarial Network) \\
GIN & Графовая изоморфная сеть (Graph Isomorphism Network) \\
GNN (ГНС) & Графовая нейронная сеть (Graph Neural Network) \\
GPU & Графический процессор (Graphics Processing Unit) \\
HIDPS (ГСОПВ) & Гибридная система обнаружения и предотвращения вторжений \\
HiPPO & Оператор полиномиальной проекции высокого порядка (High-order Polynomial Projection Operator) \\
IAT & Межпакетный интервал (Inter-Arrival Time) \\
ICS (АСУ ТП) & Промышленная система управления (Industrial Control System) \\
ICS3D & Intrusion and Counter-measure Simulation in Cloud and Edge Computing Security Dataset \\
IDPS (СОПВ) & Система обнаружения и предотвращения вторжений \\
IDS (СОВ) & Система обнаружения вторжений (Intrusion Detection System) \\
IIoT & Промышленный Интернет вещей (Industrial Internet of Things) \\
IoT & Интернет вещей (Internet of Things) \\
KL & Дивергенция Кульбака--Лейблера (Kullback--Leibler Divergence) \\
LLM & Большая языковая модель (Large Language Model) \\
LSTM & Долгая краткосрочная память (Long Short-Term Memory) \\
LWE & Обучение с ошибками (Learning With Errors) \\
М1--М7 & Обозначения семи разработанных методов (см. Рис.~\ref{fig:dissertation_architecture}) \\
MA-CPO & Мультиагентная оптимизация политики с ограничениями \\
MA-PQC-IDS & Мультиагентная постквантовая криптографическая система обнаружения вторжений \\
MambaShield & Селективная модель пространства состояний с устойчивостью к отравлению (Метод М4) \\
MC (МК) & Монте-Карло \\
MC-Dropout & Дропаут Монте-Карло (Monte Carlo Dropout) \\
MCC & Коэффициент корреляции Мэтьюза (Matthews Correlation Coefficient) \\
MITRE ATT\&CK & Состязательные тактики, техники и общие знания \\
ML (МО) & Машинное обучение (Machine Learning) \\
MLP & Многослойный перцептрон (Multi-Layer Perceptron) \\
MLSecOps & Операции безопасности машинного обучения \\
NIDPS (ССОПВ) & Сетевая система обнаружения и предотвращения вторжений \\
NIST & Национальный институт стандартов и технологий (National Institute of Standards and Technology) \\
NLP & Обработка естественного языка (Natural Language Processing) \\
NPLR & Нормальное плюс низкоранговое (Normal Plus Low-Rank) \\
Нейронное ОДУ & Нейронное обыкновенное дифференциальное уравнение (Neural Ordinary Differential Equation) \\
ОДУ & Обыкновенное дифференциальное уравнение (Ordinary Differential Equation) \\
ODE-S6 & ОДУ-дополненное селективное пространство состояний (блок CyberSecLLM) \\
OT (ОТ) & Оптимальный транспорт (Optimal Transport) \\
OWASP & Open Worldwide Application Security Project \\
PAC & Вероятно приблизительно корректный (Probably Approximately Correct) \\
PARD & Прогрессивная состязательная дистилляция робастности \\
PCAP & Захват пакетов (Packet Capture) \\
PGD & Проецированный градиентный спуск (Projected Gradient Descent) \\
PQC & Постквантовая криптография (Post-Quantum Cryptography) \\
PWA & Прогрессивное веб-приложение (Progressive Web Application) \\
Т1--Т4 & Обозначения требований: Робастность, Точность, Эффективность, Приватность \\
RAG & Генерация с дополненным извлечением (Retrieval-Augmented Generation) \\
RL & Обучение с подкреплением (Reinforcement Learning) \\
RMF & Система управления рисками (NIST AI RMF) \\
RNN (РНС) & Рекуррентная нейронная сеть (Recurrent Neural Network) \\
SCADA & Диспетчерское управление и сбор данных (Supervisory Control and Data Acquisition) \\
SDE (СДУ) & Стохастическое дифференциальное уравнение (Stochastic Differential Equation) \\
SDE-TGNN & Темпоральная графовая нейронная сеть на стохастических дифференциальных уравнениях \\
SOC & Центр операций безопасности (Security Operations Centre) \\
SSM & Модель пространства состояний (State-Space Model) \\
Стох.\ трансф. & Стохастический трансформер (Метод М5) \\
SVM & Метод опорных векторов (Support Vector Machine) \\
TABN & Темпоральная адаптивная пакетная нормализация (Temporal Adaptive Batch Normalisation) \\
TGNN & Темпоральная графовая нейронная сеть (Temporal Graph Neural Network) \\
TLS & Протокол безопасности транспортного уровня (Transport Layer Security) \\
TripleE-TGNN & Темпоральная графовая нейронная сеть с тройным вложением для многоуровневых представлений (Метод М3) \\
UC-HGP & Иерархический гауссовский процесс с калиброванной неопределённостью (Метод М6) \\
\end{longtable}
